% BANKS-SOURCE: pdf=106 print=96 kind=section
\subsection{The non-abelian Higgs phenomenon}

% BANKS-SOURCE: pdf=106 print=96 kind=prose
In a theory with gauge group $G$, the scalar fields $\Phi$ will transform in a
(generally reducible) representation $R_S$ of $G$. In this section we will call
the generators of $G$ in this representation $T^a$. In later sections we will
call them $T_S^a$. Given a point $v$ in the (linear) space of scalars, let $H$
be the stability subgroup of $v$ ($hv=v$ if $h\in H$). Its generators are
$t^i$, which are a subset of the $T^a$. We can write the general scalar field
as
\begin{equation*}
  \Phi(x)=\Omega(x)[v+\Delta(x)],
\end{equation*}
where $\Delta$ is in the subspace of field space orthogonal to all of the
vectors $T^av$. Under a gauge transformation $\Omega(x)\to V(x)\Omega(x)$,
while $\Delta(x)$ is invariant.

This parametrization of field space has a gauge ambiguity
$\Omega(x)\to\Omega(x)h(x)$, $\Delta(x)\to h^{-1}(x)\Delta(x)$, with
$h(x)\in H$. The new gauge group is isomorphic to $H$. This fact, and the close
relationship of the mathematics to that of a theory with global symmetry $G$
broken down to $H$ by the VEV $\langle\Phi\rangle=v$, accounts for the
standard terminology ``spontaneously broken gauge symmetry.'' We say that a
semi-classical expansion around the point $\Delta(x)=0$ \emph{spontaneously
breaks the gauge group $G$ to the gauge group $H$}. In fact what has happened
is that our parametrization of the $G$ gauge-invariant field space $\Delta$ has
a new $H$ gauge ambiguity.

Now define
\begin{equation*}
  B_\mu=\Omega^{-1}D_\mu\Omega\equiv a_\mu+W_\mu.
\end{equation*}
The Lie algebra of $G$ is the direct sum of the Lie algebra of $H$ and the
subspace of generators orthogonal to those in $H$. We call these the
\emph{coset generators} for the coset $G/H$. The coset generators\footnote[2]{The
coset is not generally a group, but the coset generators are a linear subspace
(not a subalgebra) of the Lie algebra.} form a representation of $H$, under the
adjoint action of $H$ on the Lie algebra of $G$. The second equality refers to
this decomposition. $a_\mu$ involves only $H$ generators, while $W_\mu$ is
composed of coset generators. All components of $B_\mu$ are invariant under
$G$ gauge transformations. Under the new $H$ gauge transformations $a_\mu$
transforms like a connection or gauge potential, while $W_\mu$ transforms
homogeneously.

Now we can rewrite the scalar field kinetic term as
\begin{equation*}
  -\frac12(D_\mu(A)\Phi)^2
  =-\frac12\left[\Omega\bigl(D_\mu(a)\Delta
  +W_\mu(v+\Delta)\bigr)\right]^2.
\end{equation*}
Using the facts that $v$ is annihilated by all generators of $H$ and that
$\Delta^TT^av=0$ for all generators of $G$, we can rewrite this as
\begin{equation*}
  -(D_\mu(A)\Phi)^2=-(D_\mu(a)\Delta)^2
  -2\bigl(W_\mu(v+\Delta)\bigr)^TD_\mu(a)\Delta
  -(W_\mu v)^2-(W_\mu\Delta)^2.
\end{equation*}

% BANKS-SOURCE: pdf=107 print=97 kind=prose
On expanding around $\Delta=0$, we find that the $W_\mu$ fields are massive
vectors, with mass matrix\footnote[3]{Actually, we are writing these formulae for
gauge potentials whose kinetic term involves a factor of the gauge coupling, so
the properly normalized mass matrix includes a factor of $g^ag^b$, where $g^a$
is the coupling of the simple factor of $G$ to which $T^a$ belongs.}
\begin{equation*}
  (\mu_V^2)^{ab}=v^TT^aT^bv.
\end{equation*}
The $a_\mu$ fields do not get a mass. The fields $\Delta$ are $G$-gauge-invariant,
and the gauge-invariant potential $V(\Phi)$, whose minima are at $\Phi=\Omega v$,
will generally give mass to all of them. The $\Delta$ fields are the physical
Higgs bosons.

So, given a non-abelian gauge theory with scalar fields, the space of
gauge-equivalent classical vacuum states is the coset space $G/H$ of fields
$\Phi=\Omega v$. The Higgs fields $\Delta$ parametrize gauge-invariant
deformations of the scalar field away from this vacuum manifold. The number of
massive gauge bosons $W_\mu$ is just the dimension of $G/H$, and there are
perturbatively massless vector potentials $a_\mu$ for the stability subgroup
$H$ of $v$.
